\begin{theorem}[corrected CDP global-vanishing criterion]
\label{thm:corrected-cdp-global}
\Verified\ Let \(X\) be a compact connected complex threefold such that
\[
 H^1(X;\mathbb Z)=H^2(X;\mathbb Z)=0,\qquad a(X)=1,
\]
and suppose its algebraic reduction is a holomorphic map
\(f:X\to C\).  If there exists \(L_0\in\operatorname{Pic}(X)\) with
\[
 H^2(X,T_X\otimes L_0)=0,
\tag{5.24}
\]
then, for \(M\) in a dense open subset of \(\operatorname{Pic}(X)\),
\[
 H^i(X,T_X\otimes M)=0\quad(i\ne1),
\tag{5.25}
\]
\[
 \chi(X,T_X\otimes M)\le0,\qquad
 \left\langle c_3(T_X),[X]\right\rangle\le0.
\tag{5.26}
\]
\end{theorem}

\begin{proof}
The exponential sequence and the two integral vanishings give
\[
 \operatorname{Pic}(X)=\operatorname{Pic}^0(X)
 \simeq H^1(X,\mathcal O_X).
\tag{5.27}
\]
Thus this Picard group is a connected irreducible complex vector group.  Any fibre divisor \(D\) of
\(f\) is nonzero and effective, and its line bundle lies in
\(\operatorname{Pic}^0(X)\) by \(H^2(X;\mathbb Z)=0\).  Applying
\cref{lem:cdp-generic-endpoints} to \(B=T_X\) shows that the degree-zero and
degree-three groups in (5.25) vanish away from two proper analytic subsets
of \(\operatorname{Pic}(X)\).

The cohomology-jump locus
\[
 \{M\in\operatorname{Pic}(X):h^2(X,T_X\otimes M)\ge1\}
\]
is analytic by upper semicontinuity.  It is proper because (5.24) gives a
point outside it.  Its complement therefore meets the two endpoint
vanishing loci in a dense open set.  This proves (5.25), and on that set
\[
 \chi(X,T_X\otimes M)=-h^1(X,T_X\otimes M)\le0.
\tag{5.28}
\]

Every \(M\in\operatorname{Pic}(X)\) has rationally trivial first Chern
class.  The hypotheses also give \(c_1(T_X)=0\) in rational cohomology and,
by Poincar\'e duality, \(H^4(X;\mathbb Q)=0\), so
\(c_2(T_X)=0\) rationally.  The degree-six term of
\(\operatorname{ch}(T_X\otimes M)\operatorname{td}(X)\) is consequently
\(\tfrac12c_3(T_X)\).  Hirzebruch--Riemann--Roch turns (5.28) into
\[
 \chi(X,T_X\otimes M)
 =\frac12\left\langle c_3(T_X),[X]\right\rangle\le0,
\]
which proves (5.26).
\end{proof}

